\section{Phase space structure of EM}

Let our spacetime be a four dimensional manifold $\Omega$, with boundary $\partial\Omega = \Sigma_-\cup\Sigma_+\cup\Gamma$. 
Our action principle is given by,

\begin{align}
    S[A] = \int\limits_\Omega\dd[4]{x}\mathcal L[\phi(x),x] = -\frac14\int\limits_\Omega\dd[4]{x} F^{\mu\nu}F_{\mu\nu}
\end{align}

The variation of the action cannot have contributions from the bulk of $\Gamma$, and also is done with compact support 
outside of $\Sigma_-\cup\Sigma_+$, that is, on $\tilde\Omega = \Omega'\cup\Gamma'$,

\begin{align}
    \int\limits_{\tilde\Omega}\dd[4]{y}\fdv{S[A]}{A_\alpha\qty(y)}\psi_\alpha\qty(y)&=\int\limits_{\tilde\Omega}\dd[4]{y}\int\limits_\Omega\dd[4]{x}\fdv{\mathcal L[x]}{A_\alpha\qty(y)}\psi_\alpha(y)\\
    \int\limits_{\tilde\Omega}\dd[4]{y}\fdv{S[A]}{A_\alpha\qty(y)}\psi_\alpha\qty(y)&=\int\limits_{\tilde\Omega}\dd[4]{y}\int\limits_\Omega\dd[4]{x}\qty[\pdv{\mathcal L[x]}{A_\alpha\qty(x)}\delta^{(4)}\qty(x-y)+\pdv{\mathcal L[x]}{\partial_\mu A_\alpha\qty(x)}\partial_\mu\delta^{(4)}\qty(x-y)]\psi_\alpha(y)\\
    \int\limits_{\tilde\Omega}\dd[4]{y}\fdv{S[A]}{A_\alpha\qty(y)}\psi_\alpha\qty(y)&=\int\limits_{\tilde\Omega}\dd[4]{y}\int\limits_\Omega\dd[4]{x}\qty[\pdv{\mathcal L[x]}{A_\alpha\qty(x)}-\partial_\mu\pdv{\mathcal L[x]}{\partial_\mu A_\alpha\qty(x)}]\delta^{(4)}\qty(x-y)\psi_\alpha(y)\nonumber\\
    &\quad\quad\quad+\int\limits_{\tilde\Omega}\dd[4]{y}\int\limits_\Omega\dd[4]{x}\partial_{x\mu}\qty[\pdv{\mathcal L[x]}{\partial_\mu A_\alpha\qty(x)}\delta^{(4)}\qty(x-y)]\psi_\alpha(y)\\
    \int\limits_{\tilde\Omega}\dd[4]{y}\fdv{S[A]}{A_\alpha\qty(y)}\psi_\alpha\qty(y)&=\int\limits_{\tilde\Omega}\dd[4]{y}\qty[\pdv{\mathcal L[y]}{A_\alpha\qty(y)}-\partial_\mu\pdv{\mathcal L[y]}{\partial_\mu A_\alpha\qty(y)}]\psi_\alpha(y)\nonumber\\
    &\quad\quad\quad+\int\limits_{\tilde\Omega}\dd[4]{y}\int\limits_{\partial\Omega}\dd[3]{S_{x\mu}}\pdv{\mathcal L[x]}{\partial_\mu A_\alpha\qty(x)}\delta^{(4)}\qty(x-y)\psi_\alpha(y)\\
    \int\limits_{\tilde\Omega}\dd[4]{y}\fdv{S[A]}{A_\alpha\qty(y)}\psi_\alpha\qty(y)&=\int\limits_{\tilde\Omega}\dd[4]{y}\qty[\pdv{\mathcal L[y]}{A_\alpha\qty(y)}-\partial_\mu\pdv{\mathcal L[y]}{\partial_\mu A_\alpha\qty(y)}]\psi_\alpha(y)\nonumber\\
    &\quad\quad\quad+\int\limits_{\partial\Omega}\dd[3]{S_{\mu}}\pdv{\mathcal L[x]}{\partial_\mu A_\alpha\qty(x)}\psi_\alpha(x)\\
    \int\limits_{\tilde\Omega}\dd[4]{y}\fdv{S[A]}{A_\alpha\qty(y)}\psi_\alpha\qty(y)&=\int\limits_{\tilde\Omega}\dd[4]{y}E^\alpha\qty[A(x);x]\psi_\alpha(y)\nonumber\\
    &\quad\quad\quad+\int\limits_{\partial\Omega}\dd[3]{S_{\mu}}\Theta^{\mu\alpha}\psi_\alpha(x)
\end{align}

With the identifications,

\begin{align}
    E^\alpha & = \partial_\mu F^{\mu\alpha}\\
    \Theta^{\mu\alpha}&=-F^{\mu\alpha}
\end{align}

Again, we follow a similar reasoning,

\begin{align}
    \int\limits_{\Omega}\dd[4]{y}\fdv{S[A]}{A_\alpha\qty(y)}\psi_{1\alpha}\qty(y)&=\int\limits_{\Omega}\dd[4]{y}E^\alpha\qty[y]\psi_{1\alpha}(y)\nonumber\\
    &\quad\quad\quad+\int\limits_{\partial\Omega}\dd[3]{S_{\mu}}\Theta^{\mu\alpha}\qty[x]\psi_{1\alpha}(x)\\
    \int\limits_{\Omega}\dd[4]{y}\int\limits_{\Omega}\dd[4]{z}\frac{\delta^2 S[A]}{\delta A_\beta\qty(z)\delta A_\alpha\qty(y)}\psi_{1\alpha}\qty(y)\psi_{2\beta}\qty(z)&=\int\limits_{\Omega}\dd[4]{y}\int\limits_{\Omega}\dd[4]{z}\fdv{E^\alpha\qty[y]}{A_\beta(z)}\psi_{1\alpha}(y)\psi_{2\beta}\qty(z)\nonumber\\
    &\quad\quad\quad+\int\limits_{\partial\Omega}\dd[3]{S_{\mu}}\int\limits_{\Omega}\dd[4]{z}\fdv{\Theta^{\mu\alpha}\qty[x]}{A_\beta\qty(z)}\psi_{1\alpha}(x)\psi_{2\beta}\qty(z)\\
\end{align}

Antisymmetrize in $1\leftrightarrow2$,
\begin{align}
    0&=\int\limits_{\Omega}\dd[4]{y}\int\limits_{\Omega}\dd[4]{z}\qty(\fdv{E^\alpha\qty[y]}{A_\beta(z)}-\fdv{E^\beta\qty[z]}{ A_\alpha(y)})\psi_{1\alpha}(y)\psi_{2\beta}\qty(z)\nonumber\\
    &\quad\quad\quad+\int\limits_{\partial\Omega}\dd[3]{S_{\mu}}\int\limits_{\Omega}\dd[4]{z}\qty(\fdv{\Theta^{\mu\alpha}\qty[x]}{A_\beta\qty(z)}-\fdv{\Theta^{\mu\beta}\qty[z]}{A_\alpha\qty(x)})\psi_{1\alpha}(x)\psi_{2\beta}\qty(z)
\end{align}

As long as the equation of motion have just even order of derivatives with constant coefficients,
\begin{align}
    0&=-\int\limits_{\partial\Omega}\dd[3]{S_{\mu}}\int\limits_{\Omega}\dd[4]{z}\qty(g^{\alpha\beta}\partial^\mu_x\delta^{(4)}\qty(x-z)-g^{\mu\beta}\partial^\alpha_x\delta^{(4)}\qty(x-z)-g^{\beta\alpha}\partial^\mu_z\delta^{(4)}\qty(z-x)+g^{\mu\alpha}\partial^\beta_z\delta^{(4)}\qty(z-x))\psi_{1\alpha}(x)\psi_{2\beta}\qty(z)\\
    0&=-\int\limits_{\partial\Omega}\dd[3]{S_{\mu}}\qty(\psi_{1\alpha}g^{\alpha\beta}\partial^\mu\psi_{2\beta}-\psi_{1\alpha}g^{\mu\beta}\partial^\alpha\psi_{2\beta}-\psi_{2\beta}g^{\beta\alpha}\partial^\mu\psi_{1\alpha}+\psi_{2\beta}g^{\mu\alpha}\partial^\beta\psi_{1\alpha})
\end{align}

Suppose that for a given solution of the equations of motion $A$, it's true that $A+\psi_i$ are also solutions of the 
equations of motion, then,
\begin{align}
    \int\limits_{\Sigma_+}\dd[3]{S_{\mu}}\qty(\psi_{1\alpha}\qty(\partial^\mu\psi_{2}^\alpha-\partial^\alpha\psi_{2}^\mu)-\psi_{2\alpha}\qty(\partial^\mu\psi_{1}^\alpha-\partial^\alpha\psi_{1}^\mu))=\int\limits_{\Sigma_-}\dd[3]{S_{\mu}}\qty(\psi_{1\alpha}\qty(\partial^\mu\psi_{2}^\alpha-\partial^\alpha\psi_{2}^\mu)-\psi_{2\alpha}\qty(\partial^\mu\psi_{1}^\alpha-\partial^\alpha\psi_{1}^\mu))
\end{align}

In other words, the functional,
\begin{align}
    \Omega\qty(\psi_1,\psi_2) &= \int\limits_{\Sigma_t}\dd[3]{\vb x}\qty(\psi_{1i}\qty(\dot\psi_2^i-\partial^i\psi_{20})-\psi_{2i}\qty(\dot\psi_1^i-\partial^i\psi_{10}))
\end{align}
defined in the tangent space of the phase space of solutions, is independent of the time slice used to compute it. The main 
problem here is that this functional is degenerated in the variables $\psi_\alpha$, this means that they aren't a good parametrization 
of the phase space. To see that are degeneracies here, 
\begin{align}
    \Omega\qty(\psi_1+\partial\lambda,\psi_2) &= \Omega\qty(\psi_1,\psi_2)+\int\limits_{\Sigma_t}\dd[3]{\vb x}\qty(\partial_i\lambda\qty(\dot\psi_2^i-\partial^i\psi_{20})-\psi_{2i}\qty(\partial^i\dot\lambda-\partial^i\dot\lambda))\\
    \Omega\qty(\psi_1+\partial\lambda,\psi_2) &= \Omega\qty(\psi_1,\psi_2)+\int\limits_{\partial\Sigma_t}\dd[2]{S_i}\lambda\qty(\dot\psi_2^i-\partial^i\psi_{20})-\int\limits_{\Sigma_t}\dd[3]{\vb x}\lambda\partial_i\qty(\dot\psi_2^i-\partial^i\psi_{20})\\
    \Omega\qty(\psi_1+\partial\lambda,\psi_2) &= \Omega\qty(\psi_1,\psi_2)+\int\limits_{\partial\Sigma_t}\dd[2]{S_i}\lambda\qty(\dot\psi_2^i-\partial^i\psi_{20})
\end{align}
as long as $\lambda$ is well behaved in the spatial boundary, the last integral is zero, and thus we have that the $\Omega$ is 
degenerated in this space, also, we can shift as we please the canonical momentum $\dot\psi_0$, with no impact on $\Omega$. In 
total, we should have just two coordinates and conjugated momenta, the usual choice is,








Taken 
into account that the phase space needs two variables for each degree of freedom, they are chosen as $\phi(\vb x)$ and $\pi\qty(\vb x)=\dot\phi\qty(\vb x)$, 
in this way, our functional can be rewritten as a phase space 2-form,

\begin{align}
    \Omega &= \int\limits_{\Sigma_t}\dd[3]{\vb x}\delta\phi\qty(\vb x)\wedge\delta\pi\qty(\vb x)
\end{align}

The main point is that this 2-form is non degenerated and can be inverted as,

\begin{align}
    \acomm{\cdot}{\cdot} &= \int\limits_{\Sigma_t}\dd[3]{\vb x}\qty(\fdv{}{\phi\qty(\vb x)}\fdv{}{\pi\qty(\vb x)}-\fdv{}{\pi\qty(\vb x)}\fdv{}{\phi\qty(\vb x)})
\end{align}

Which provides the fundamental Poisson bracket,
\begin{align}
    \acomm{\phi\qty(t,\vb x)}{\pi\qty(t,\vb y)}&=\delta^{(3)}\qty(\vb x-\vb y) 
\end{align}
% \int\limits_{\Omega}\dd[4]{y}\int\limits_{\Omega}\dd[4]{z}\frac{\delta^2 S[\phi]}{\delta\phi\qty(z)\delta\phi\qty(y)}\psi_1\qty(y)\psi_2\qty(z)&=\int\limits_{\Omega}\dd[4]{y}\int\limits_{\Omega}\dd[4]{z}\qty(\Box_y-m^2)\delta^{(4)}\qty(y-z)\psi_1(y)\psi_2\qty(z)\nonumber\\
    % &\quad\quad\quad-\int\limits_{\partial\Omega}\dd[3]{S_{\alpha}}\int\limits_{\Omega}\dd[4]{z}\partial^\alpha_x\delta^{(4)}\qty(x-z)\psi_1(x)\psi_2\qty(z)\\
    % \int\limits_{\Omega}\dd[4]{y}\int\limits_{\Omega}\dd[4]{z}\frac{\delta^2 S[\phi]}{\delta\phi\qty(z)\delta\phi\qty(y)}\psi_1\qty(y)\psi_2\qty(z)&=\int\limits_{\Omega}\dd[4]{y}\int\limits_{\Omega}\dd[4]{z}\psi_1(y)\fdv{\psi_2\qty(y)}{\phi(z)}\qty(\Box_z-m^2)\phi(z)\nonumber\\
    % &\quad\quad\quad-\int\limits_{\partial\Omega}\dd[3]{S_{\alpha}}\int\limits_{\Omega}\dd[4]{z}\psi_1(x)\fdv{\psi_2\qty(x)}{\phi(z)}\partial^\alpha_z\phi\qty(z)

% Evaluate at a solution of the equations of motion,

% \begin{align}
%     \int\limits_{\Omega}\dd[4]{y}\int\limits_{\Omega}\dd[4]{z}\frac{\delta^2 S[\phi]}{\delta\phi\qty(z)\delta\phi\qty(y)}\psi_1\qty(y)\psi_2\qty(z)&=-\int\limits_{\Sigma_+\cup\Sigma_-}\dd[3]{S_{\alpha}}\int\limits_{\Omega}\dd[4]{z}\psi_1(x)\fdv{\psi_2\qty(x)}{\phi(z)}\partial^\alpha_z\phi\qty(z)
% \end{align}

% This should be symmetric in $1\leftrightarrow2$, thus,

% \begin{align}
%     0&=-\int\limits_{\Sigma_+\cup\Sigma_-}\dd[3]{S_{\alpha}}\int\limits_{\Omega}\dd[4]{z}\psi_1(x)\fdv{\psi_2\qty(x)}{\phi(z)}\partial^\alpha_z\phi\qty(z)+\int\limits_{\Sigma_+\cup\Sigma_-}\dd[3]{S_{\alpha}}\int\limits_{\Omega}\dd[4]{z}\psi_2(x)\fdv{\psi_1\qty(x)}{\phi(z)}\partial^\alpha_z\phi\qty(z)
% \end{align}

% \begin{align}
%     &\int\limits_{\Sigma_+}\dd[3]{S_{\alpha}}\int\limits_{\Omega}\dd[4]{z}\partial^\alpha_z\phi\qty(z)\qty[\psi_1(x)\fdv{\psi_2\qty(x)}{\phi(z)}-\psi_2(x)\fdv{\psi_1\qty(x)}{\phi(z)}]=\int\limits_{\Sigma_-}\dd[3]{S_{\alpha}}\int\limits_{\Omega}\dd[4]{z}\partial^\alpha_z\phi\qty(z)\qty[\psi_1(x)\fdv{\psi_2\qty(x)}{\phi(z)}-\psi_2(x)\fdv{\psi_1\qty(x)}{\phi(z)}]\\
%     &\int\limits_{\Sigma_+}\dd[3]{S_{\alpha}}\qty[\psi_1(x)\partial^\alpha\psi_2\qty(x)-\psi_2(x)\partial^\alpha\psi_1\qty(x)]=\int\limits_{\Sigma_-}\dd[3]{S_{\alpha}}\qty[\psi_1(x)\partial^\alpha\psi_2\qty(x)-\psi_2(x)\partial^\alpha\psi_1\qty(x)]
% \end{align}